\TieudeT{PHONG TỎA VẬT LÍ 12 - VẬT LÍ NHIỆT}
\subsubsection*{PHẦN I. Thí sinh trả lời câu hỏi từ câu 1 đến câu 18. Mỗi câu hỏi thí sinh chỉ chọn một phương án}
\setcounter{ex}{0}
\begin{ex}
Nội năng của một vật phụ thuộc vào
\choice
{khối lượng và thể tích của vật}
{\True nhiệt độ và thể tích của vật}
{nhiệt độ và khối lượng riêng của vật}
{khối lượng và khối lượng riêng của vật}
\loigiai{}
\end{ex}

\begin{ex}
Quá trình làm thay đổi nội năng của vật bằng cách cho nó tiếp xúc với vật khác khi
\choice
{nhiệt độ của chúng bằng nhau gọi là sự trao đổi công}
{có sự chênh lệch nhiệt độ giữa chúng gọi là sự nhận công}
{\True có sự chênh lệch nhiệt độ giữa chúng gọi là sự truyền nhiệt}
{nhiệt độ của chúng bằng nhau gọi là sự truyền nhiệt}
\loigiai{}
\end{ex}

\begin{ex}
Phát biểu nào sau đây \textbf{không} đúng với mô hình động học phân tử?
\choice
{Các chất được cấu tạo từ các hạt riêng biệt là phân tử}
{Các phân tử chuyển động không ngừng}
{\True Tốc độ chuyển động của các phân tử cấu tạo nên vật càng lớn thì thể tích của vật càng lớn}
{Giữa các phân tử có lực tương tác gọi là lực tương tác phân tử}
\loigiai{}
\end{ex}

\begin{ex}
Với mô hình động học phân tử, sự khác biệt về độ lớn của lực tương tác giữa các phân tử trong chất rắn, chất lỏng, chất khí dẫn đến sự
\choice
{đồng nhất về cấu trúc của chúng}
{\True khác biệt về cấu trúc của chúng}
{khác biệt về khối lượng của chúng}
{đồng nhất về khối lượng của chúng}
\loigiai{}
\end{ex}

\begin{ex}
Phát biểu nào sau đây là \textbf{sai}?
\choice
{Khoảng cách giữa các phân tử càng lớn thì lực tương tác giữa chúng càng yếu}
{Các phân tử sắp xếp càng có trật tự thì lực tương tác giữa chúng càng mạnh}
{\True Lực tương tác giữa các phân tử không thể là lực đẩy}
{Khi khoảng cách giữa các phân tử đủ lớn thì lực tương tác giữa các phân tử bằng 0}
\loigiai{}
\end{ex}

\begin{ex}
\immini[thm]{Hình bên là đồ thị sự thay đổi nhiệt độ của vật rắn kết tinh khi được làm nóng chảy. Trong khoảng thời gian từ $t_{a}$ đến $t_{b}$ thì
\choice
{vật rắn không nhận năng lượng}
{nhiệt độ của vật rắn tăng}
{nhiệt độ của vật rắn giảm}
{\True vật rắn đang nóng chảy}}{\begin{tikzpicture}[draw=Mapcolor, very thick, scale = 1, line join=round, line cap=round, >=stealth]
 \path
 (0,0)coordinate(O)
 (5,0)coordinate(x)
 (0,4)coordinate(y);
 \draw [->](O)--(x);
 \draw [->](O)--(y);
 \node at (5.5,-0.5) {Thời gian};
 \node at (0,4.5) {Nhiệt độ};
 \draw[dashed] (0,2.5) -- (1.75,2.5);
 \draw[dashed] (0.5,1.5) -- (0.5,0);
 \draw[dashed] (1.75,2.5) -- (1.75,0);
 \draw[dashed] (3.5,2.5) -- (3.5,0);
 \draw[dashed] (4.25,3.75) -- (4.25,0);
 \draw[very thick,red] (0.5,1.5) .. controls (1.5,1.75) and (0.75,2.25) .. (1.75,2.5);
 \draw[very thick,red] (1.75,2.5) -- (3.5,2.5);
 \draw[very thick,red] (3.5,2.5) .. controls (4,2.75) and (4.25,3) .. (4.25,3.75);
 \node[draw=white,align=left] at (-1,2.5) {Nhiệt độ\\ nóng chảy};
 \node[below] at (0.5,0) {$t_0$};
 \node[below] at (1.75,0) {$t_a$};
 \node[below] at (3.5,0) {$t_b$};
 \node[below] at (4.25,0) {$t_1$};
 \node[below left] at (0,0) {$0$};
 \end{tikzpicture}}

\loigiai{}
\end{ex}

\begin{ex}
Chất nào sau đây ở thể khí khi ở điều kiện thường?
\choice
{Muối ăn}
{Đường ăn}
{Giấm}
{\True Carbon dioxide}
\loigiai{}
\end{ex}

\begin{ex}
Trong các chất sau: muối ăn, nhôm, nước đá và cao su, chất rắn vô định hình là
\choice
{muối ăn}
{nhôm}
{nước đá}
{\True cao su}
\loigiai{}
\end{ex}

\begin{ex}
Sự ngưng tụ là sự chuyển từ
\choice
{thể hơi sang thể rắn}
{thể lỏng sang thể hơi}
{\True thể hơi sang thể lỏng}
{thể hơi sang thể khí}
\loigiai{}
\end{ex}

\begin{ex}
Đơn vị của nội năng là
\choice
{$^{\circ}C$}
{$K$}
{\True $J$}
{$Pa$}
\loigiai{}
\end{ex}

\begin{ex}
Khi làm nóng liên tục vật rắn vô định hình, vật rắn mềm đi và chuyển dần sang thể lỏng một cách liên tục. Trong quá trình này nhiệt độ của vật ...(1).... Do đó, vật rắn vô định hình ...(2).... Điền vào chỗ trống các cụm từ thích hợp.
\choice
{\True (1) tăng lên liên tục; (2) không có nhiệt độ nóng chảy xác định}
{(1) giữ ổn định; (2) không có nhiệt độ nóng chảy xác định}
{(1) giữ ổn định; (2) có nhiệt độ nóng chảy xác định được}
{(1) tăng lên liên tục; (2) có nhiệt độ nóng chảy xác định được}
\loigiai{}
\end{ex}

\begin{ex}
Phát biểu nào sau đây nói về nhiệt lượng là \textbf{không} đúng?
\choice
{Nhiệt lượng là số đo độ biến thiên nội năng của vật trong quá trình truyền nhiệt}
{\True Một vật lúc nào cũng có nội năng, do đó lúc nào cũng có nhiệt lượng}
{Đơn vị của nhiệt lượng cũng là đơn vị của nội năng}
{Nhiệt lượng không phải là nội năng}
\loigiai{}
\end{ex}

\begin{ex}
Lực tương tác giữa các phân tử chất rắn ... (1)... nên giữ được các phân tử ở các vị trí cân bằng và mỗi phân tử ...(2).... Điền vào chỗ trống các cụm từ thích hợp.
\choice
{(1) là lực hút; (2) dao động xung quanh vị trí cân bằng có thể di chuyển được}
{(1) rất mạnh; (2) đứng yên tại vị trí cân bằng này}
{(1) là lực hút; (2) chỉ có thể dao động xung quanh vị trí cân bằng xác định này}
{\True (1) rất mạnh; (2) chỉ có thể dao động xung quanh vị trí cân bằng xác định này}
\loigiai{}
\end{ex}

\begin{ex}
Phát biểu nào sau đây là \textbf{sai}?
\choice
{Chất rắn vô định hình không có nhiệt độ nóng chảy xác định}
{Khi bị làm nóng thì chất rắn vô định hình mềm dần cho đến khi trở thành lỏng}
{Trong quá trình hóa lỏng nhiệt độ của chất rắn vô định hình tăng liên tục}
{\True Chất rắn vô định hình có cấu trúc tinh thể}
\loigiai{}
\end{ex}

\begin{ex}
Khi bắt đầu đun, nhiệt độ của vật rắn kết tinh tăng dần. Đến nhiệt độ xác định, sự nóng chảy diễn ra, vật chuyển từ thể rắn sang thể lỏng và nhiệt độ ...(1)... dù tiếp tục đun. Sau khi toàn bộ vật chuyển sang thể lỏng, nhiệt độ của chất lỏng ...(2)... khi tiếp tục đun. Chỗ trống (1) và (2) lần lượt là
\choice
{"giảm xuống" và "giữ giá trị ổn định"}
{"không tăng" và "giảm xuống"}
{"giảm xuống" và "tiếp tục tăng lên"}
{\True "không tăng" và "tiếp tục tăng lên"}
\loigiai{}
\end{ex}

\begin{ex}
Một số phân tử ở gần mặt thoáng chất lỏng, chuyển động hướng ra ngoài, có ..(1).. đủ lớn thắng được lực tương tác giữa các phân tử thì có thể thoát ra ngoài khối chất lỏng. Như vậy, có thể nói sự bay hơi là sự hóa hơi xảy ra ở..(2).. của khối chất lỏng. Điền vào chỗ trống các cụm từ thích hợp.
\choice
{\True (1) động năng; (2) mặt thoáng}
{(1) thế năng; (2) mặt thoáng}
{(1) động năng; (2) trong lòng}
{(1) thế năng; (2) trong lòng}
\loigiai{}
\end{ex}

\begin{ex}
Người ta thực hiện công $1000\ J$ để nén khí trong xilanh. Khí truyền ra bên ngoài nhiệt lượng $600\ J$. Độ biến thiên nội năng của khí là
\choice
{$-1000\ J$}
{\True $400\ J$}
{$300\ J$}
{$-400\ J$}
\loigiai{}
\end{ex}

\begin{ex}
Khi truyền nhiệt lượng $6.10^{6}\ J$ cho khí trong một xilanh hình trụ thì khí nở ra đẩy pit-tông lên làm cho thể tích của khí tăng lên $0,5\ m^{3}$. Biết áp suất của khí là $8.10^{6}\ N/m^{2}$ và coi áp suất này không đổi trong quá trình khí thực hiện công. Độ biến thiên nội năng của khí là
\choice
{\True $2.10^{6}\ J$}
{$10^{5}\ J$}
{$-2.10^{6}\ J$}
{$6.10^{6}\ J$}
\loigiai{}
\end{ex}
\subsubsection*{PHẦN 2. Thí sinh trả lời từ câu 1 đến câu 4. Trong mỗi ý a), b), c), d) ở mỗi câu, thí sinh chọn đúng hoặc sai.}
\setcounter{ex}{0}
\begin{ex}
	\immini[thm]{Khi tiến hành đun một khối nước đá, một học sinh ghi lại được đồ thị sự phụ thuộc của nhiệt độ theo thời gian (từ lúc bắt đầu đun $\tau=0$) như hình bên.}{\begin{tikzpicture}[draw=Mapcolor, very thick, scale=1.2,thick,>=stealth']
			\tkzDefPoints{0/0/o, 6/0/x, 0/2.5/y, 1.5/0/a, 3/1.5/b, 5/1.5/c, 0/1.5/100, 3/0/b',5/0/c'}
			\tkzDrawSegments[very thick,->](o,x o,y)
			\tkzDrawSegments[dashed](100,c c,c' b,b')
			\tkzLabelPoint[left](o){$O$}
			\tkzLabelPoint[below](x){$\tau$ $(s)$}
			\tkzLabelPoint[above](y){$t$ ($^\circ C$)}
			\tkzDrawSegments[very thick, red](o,a a,b b,c)
			\tkzDrawPoints[size=4,fill=white](a,b,c,o,100)
			\tkzLabelPoint[above left](a){$A$}
			\tkzLabelPoint[left](100){$100$}
			\tkzLabelPoint[above](b){$B$}
			\tkzLabelPoint[above](c){$C$}
	\end{tikzpicture}}
	\choiceTF[t]
{\True Đồ thị hình bên mô tả quá trình nước nhận nhiệt lượng để tăng nhiệt độ và chuyển thể}
{Trên đoạn $OA$, khối nước đá không tăng nhiệt độ vì vậy nó không nhận nhiệt lượng từ nguồn nhiệt dùng để đun nước}
{Trên đoạn $AB$, xảy ra quá trình tan chảy của nước đá}
{Trên đoạn $BC$ là giai đoạn nước ở thể hơi}
	
	\loigiai{
		\begin{itemchoice}
			\itemch 
			\itemch Khối nước đá vẫn nhận nhiệt cung cấp cho quá trình nóng chảy.
			\itemch Giai đoạn này nước đã ở thể lỏng và đang tăng nhiệt độ.
			\itemch Giai đoạn này nước đang sôi nên tồn tại ở cả thể lỏng và khí.
		\end{itemchoice}	
	}
\end{ex}

\begin{ex}
Đồ thị thực nghiệm như hình bên biểu diễn sự thay đổi của nhiệt độ theo thời gian trong quá trình chuyển thể của benzene. Cho biết ở $12^\circ C$, benzene ở thể lỏng.
\begin{center}
\begin{tikzpicture}[draw=Mapcolor, very thick, scale = 0.8, line join=round, line cap=round, >=stealth,line width=1.5]
 \draw[->] (0,-3)--(0,4) node[above] {$t\ \left(^{\circ}C\right)$};
 \draw[->] (0,0) --(7,0) node[right] {$\tau\ \left(\text{phút}\right)$};
 \draw[dashed,line width=0.5,ystep=0.25,xstep=0] (0,-2.75) grid (6.5,3.5);
 \draw (2.2,3.5) node [above] {{\scriptsize 2 phút 15 s}} -- (2.2,-2.75) (4.2,3.5) node [above] {{\scriptsize 4 phút 07 s}} -- (4.2,-2.75);
 \foreach \x in {1,...,6}{
 \draw[dashed,blue,line width=1] (\x,3.5) -- (\x,-2.75);
 \draw[fill=black] (\x,0) circle (2pt);
 \draw (\x,0) node [below left] {\x};}
 \foreach \y in {-8,-4,0,4,8,12}{
 \draw[fill=black] (0,{\y/4}) circle (2pt);
 \draw (0,{\y/4}) node [left] {\y};}
 \foreach \x/\y in {0.8/3.35,1.15/2.35,1.85/1.75,2.2/1.5,3.3/1.5,4.2/1.5,4.4/0.5,4.85/0.25,5/0,5.15/-0.6,5.4/-1.5,6.3/-2.75}{\fill[blue] (\x,\y) circle (2pt);}
 \end{tikzpicture}
\end{center}
\choiceTFt
{Đây là quá trình hóa hơi của benzene}
{Trong giai đoạn 2 benzene ở thể khí (hơi)}
{\True Nhiệt độ xảy ra sự chuyển thể là $6^\circ C$}
{\True Từ phút thứ 1 đến phút thứ 2, benzene ở thể lỏng và hơi}
\end{ex}

\begin{ex}
\immini[thm]{Dùng tay thực hiện công $(A)$ ấn pit-tông để nén khí, đồng thời nung nóng khí trong xilanh bằng ngọn lửa đèn cồn. Theo định luật 1 của nhiệt động lực học.
\choiceTFt
{Công $A<0$ vì khí bị nén}
{Nhiệt lượng $Q<0$ vì khí bị nung nóng}
{\True Nội năng của khối khí tăng lên}
{Độ biến thiên nội năng của khí là $\Delta U=A-Q$}}{\begin{tikzpicture}[draw=Mapcolor, very thick, scale = 0.6, line join=round, line cap=round, >=stealth]% ĐÈN CỒN
\fill[blue!60] (-1,1) -- (1,1) -- (0.5,0) -- (-0.5,0) -- (-1,1);% CỒN
\draw[-,thick, line width = 2] (0,1.5) -- (-0.25,1.5) -- (-0.5,1) -- (-1,1) -- (-0.5,0) -- (0.5,0) -- (1,1) -- (0.5,1) -- (0.25,1.5) -- (0,1.5);% BÌNH ĐÈN CỒN
\draw[line width = 2.5] (0,1.65)..controls+(-90:0.5)and+(90:0.5)..(0,1.25)
..controls+(-35:1)and+(135:0.2)..(-0.35,0.5)..controls+(-35:0.2)and+(175:0.4)..(0.35,0.1);% DÂY
\fill[red] (0,1.65)..controls+(170:0.4)and+(-120:0.2)..(0,2.5)
..controls+(-100:0.2)and+(20:0.5)..(0,1.65);% LỬA
\begin{scope}[rotate=90,xshift=3.9cm,yshift=-3cm]
\draw[xshift = -1cm,yshift = 2.2cm,-,thick, line width = 6] (2.2,0.75) -- (4,0.75) -- (4,1) -- (4,0.5);
\fill[xshift = -1cm,yshift = 2.5cm,black!70](2,-0.5) rectangle (2.4,1.5);
\draw[xshift = -1cm,yshift = 2.5cm,-,thick, line width = 3,rounded corners=2pt] (3,-0.5) -- (0,-0.5) -- (0,1.5) -- (3,1.5);
\fill[xshift = -1cm,yshift = 2.5cm,pattern=crosshatch dots,-,thick, line width = 2,rounded corners=4pt] (2,-0.5) -- (0,-0.5) -- (0,1.5) -- (2,1.5);
\end{scope}
\begin{scope}[yshift=7cm]
\draw[line width = 1] (1.1,1.1) -- (0.6,1.1)..controls+(-110:0.1)and+(0:0.2)..(0.4,0.7) -- (-0.25,0.6)..controls+(-160:0.5)and+(70:0.2)..(-1.45,0.1)..controls+(-90:0.1)and+(170:0.05)..(-1.25,0.05);
\draw[line width = 1] (0.3,0.35)..controls+(-180:1.2)and+(-165:1)..(0.2,0.1)..controls+(5:1.1)and+(-110:0.6)..(1.1,1.1);
\draw[line width = 1] (-0.25,0.39)..controls+(-160:0.2)and+(60:0.2)..(-1.25,0.075)..controls+(-90:0.2)and+(170:0.2)..(-0.52,0.07);
\draw[line width = 1] (-0.3,0.28)..controls+(-40:0.1)and+(60:0.1)..(-0.3,0.05);
\draw[line width = 1] (-0.3,0.28)..controls+(-40:0.1)and+(60:0.1)..(-0.3,0.05);
\end{scope}
\end{tikzpicture}}

\end{ex}

\begin{ex}
	Một quả bóng có khối lượng $200\ g$ rơi từ độ cao $15\ m$ xuống sân và nảy lên đến độ cao $10\ m$. Sau khi va chạm với mặt sân, thể tích quả bóng không thay đổi. Bỏ qua ma sát và lực cản của không khí. Lấy $g=10\ m/s^2$.
	\choiceTFt
	{Ở độ cao $10\ m$, quả bóng không có nội năng}
	{Trong quá trình rơi xuống, nội năng của quả bóng giảm vì thế năng của quả bóng giảm}
	{Sau khi va chạm với mặt đất, nội năng của quả bóng tăng lên $20\ J$}
	{\True Sự thay đổi nội năng của quả bóng là thay đổi động năng trung bình của các phân tử mà không có thay đổi thế năng của các phân tử}
\end{ex}

\subsubsection*{PHẦN 3. Thí sinh trả lời từ câu 1 đến câu 6.}
\setcounter{ex}{0}
\begin{ex}
Người ta thực hiện công $150\ kJ$ để nén một khối khí, đồng thời cung cấp cho khối khí này một nhiệt lượng $20\ kJ$. Biết khí không tỏa nhiệt ra môi trường bên ngoài. Độ biến thiên nội năng của khối khí bằng bao nhiêu kilô jun?
\SA[oly]{170}
\end{ex}

\begin{ex}
Người ta thực hiện công $1000\ J$ để nén khí trong một xilanh. Biết khí truyền ra môi trường xung quanh nhiệt lượng $400\ J$. Độ biến thiên nội năng của khối khí bằng bao nhiêu jun?
\SA[oly]{600}
\end{ex}

\begin{ex}
Khi truyền nhiệt lượng cho khối khí trong một xilanh hình trụ thì khí dãn nở đẩy pit-tông làm thể tích của khối khí tăng thêm $7\ \text{lít}$. Biết áp suất của khối khí là $3.10^{5}\ N/m^{2}$ và không đổi trong quá trình khí dãn nở. Công mà khối khí thực hiện có độ lớn bằng bao nhiêu jun?
\SA[oly]{2100}
\end{ex}

\begin{ex}
Một khối lượng khí lí tưởng ở áp suất $3.10^{5}\ Pa$ ($1\ Pa = 1\ N/m^{2}$) có thể tích $8\ \text{lít}$. Sau khi nhận được nhiệt lượng $1000\ J$, khí dãn ra và có thể tích $10\ \text{lít}$. Coi áp suất của khối khí không đổi trong quá trình dãn nở và khí không tỏa nhiệt ra môi trường bên ngoài. Độ biến thiên nội năng của khối khí bằng bao nhiêu jun?
\SA[oly]{400}
\end{ex}

\begin{ex}
Một viên đạn chì khối lượng $m = 10\ g$ có vận tốc giảm từ $v_{1} = 400\ m/s$ xuống $v_{2} = 300\ m/s$ khi xuyên qua một tấm ván. Biết chỉ có $80\%$ công mà viên đạn sinh ra khi xuyên qua tấm ván chuyển thành nhiệt. Nhiệt lượng tỏa ra của viên đạn trong quá trình đó là bao nhiêu jun?
\SA[oly]{280}
\end{ex}

\begin{ex}
Một bình đun nước nóng bằng điện có công suất $9\ kW$. Nước ở $15^{\circ}C$ được làm nóng khi đi qua buồng đốt của bình. Nước chảy qua buồng đốt với lưu lượng $58.10^{-3}\ kg/s$. Cho biết để làm cho $1\ kg$ nước tăng thêm $1^{\circ}C$ thì cần năng lượng là $4180\ J$. Nhiệt độ của nước khi ra khỏi buồng đốt bằng bao nhiêu $^{\circ}C$ (làm tròn kết quả đến hàng phần mười)?
\SA[oly]{52,1}
\end{ex}


